% !PRINTED 400
\weberpage{400}
% Sechzehnter Abschnitt. \hfill \S\ 110. % running head / page furniture
%<<CONT

die fundamentale Einheit mit positiver Norm, also $T$, $U$ die
kleinste positive Lösung von
\weq{3}{T^2 - U^2 D = +4}
ist.

Bei positiver Diskriminante gibt es unendliche Scharen assoziierter
Zahlen, und es kommt darauf an, durch eine passende Bestimmung von
jeder dieser Scharen eine bestimmte auszusondern. Dazu bieten uns
die allgemeinen Betrachtungen von Bd.~II, \S\ 195 die Hilfsmittel,
die wir jetzt auf unseren speziellen Fall anwenden.

Es sei nach \S\ 98
\[
\lambda = a\, t_1 + \frac{b+\sqrt{D}}{2}\, t_2
\]
eine Basisform in der Ordnung $[Q]$ eines Ideals $\mathfrak{a}$, das wir
\gesperrt{primär} voraussetzen wollen und
\weq{4}{\begin{aligned}
y &= a\, x_1 + \frac{b+\sqrt{D}}{2}\, x_2 \\
y' &= a\, x_1 + \frac{b-\sqrt{D}}{2}\, x_2
\end{aligned}}
mit ganzen rationalen $x_1$, $x_2$ eine durch $\mathfrak{a}$ teilbare
ganze Zahl $y$ in $[Q]$ mit positiver Norm nebst ihrer Konjugierten
$y'$, und da wir $\mathfrak{a}$ als primär vorausgesetzt haben, so ist
\weq{5}{a = N(\mathfrak{a}).}
Es sei ferner
\weq{6}{\begin{aligned}
\varepsilon &= \frac{T+U\sqrt{D}}{2}, \\
\varepsilon' &= \frac{T-U\sqrt{D}}{2}
\end{aligned}}
die fundamentale Einheit in $O$ mit ihrer Konjugierten, die beide
\gesperrt{positiv sind} und positive Norm haben.

Wir bestimmen die Zahlen $\xi_1$ und $\xi_2$ aus den linearen
Gleichungen:
\weq{7}{\begin{aligned}
\xi_1 \log \varepsilon + \xi_2 &= \log |y|, \\
\xi_1 \log \varepsilon' + \xi_2 &= \log |y'|, \\
\xi_1 \log \frac{\varepsilon}{\varepsilon'} &= \log \left| \frac{y}{y'} \right|, \\
\xi_2 &= \tfrac{1}{2} \log |y y'|,
\end{aligned}}
worin $|y|$ den absoluten Wert von $y$ bedeutet.
